\documentclass{zhvt-classic}

% 使用本地源流明体（GenRyuMin2 TW）按文件路径加载，不依赖系统字体名
% 版框（内框／界栏／书口侧栏／鱼尾）改用与版心印章同一印泥棕色：
%   印章图片 yshs.png 实色 = RGB(135,68,52) = HTML 874434
\usepackage{xcolor}
\usepackage{tikz}
% sealbrown 双注册：xcolor 供 TikZ/pgf，L3 供 cls 版框
\definecolor{sealbrown}{HTML}{874434}
\ExplSyntaxOn
\color_set:nnn { sealbrown } { HTML } { 874434 }
\ExplSyntaxOff

\zhvtset{
  grid color       = sealbrown,
  main font        = GenRyuMin2TW-SB,
  main font path   = C:/Users/Administrator/texmf/fonts/opentype/public/GenRyuMin2TW/,
jiazhu font      = GenRyuMin2TW-M,
  jiazhu font path = C:/Users/Administrator/texmf/fonts/opentype/public/GenRyuMin2TW/,
  % 页码框容纳 4 字：正常页「九十七」+前置空 1 格=4 行恰好；
  % 100+ 页「一百零一」共 5 行超出，末字「一/二/三」被框裁掉（2026-10-04 实测），故放宽到 5
  shukou page number chars = 5,
  % 书口书名框高从默认 6 字减到 4 字，把省出的 2 字高拨给 chapter 框，
  % 使 6 字书名+卷名放进 9 字高的 chapter 框后不再压线（2026-10-05）
  shukou title chars = 4,
}

\title{史记}[司马相如列传]
\author{汉·司马迁}
\maker{闲叟漫记·重製}


% ===== 索隐述赞样式（2026-11-08 四次定稿 · 愚大人亲手改定参数）=====
% ⚠ 竖排前提：自定义字体家族必须带 RawFeature={vertical:+vert}；
%   simhei.ttf 加 +vhal 会整体加载失败静默回退，只写 +vert；simfang 同理。
\newCJKfontfamily{\shuzanfont}[Path=C:/Windows/Fonts/,RawFeature={vertical:+vert}]{simhei.ttf}
\newCJKfontfamily{\shuzanbody}[Path=C:/Windows/Fonts/,RawFeature={vertical:+vert}]{simfang.ttf}
% 题头：152x30pt 圆角 6pt 棕底白字；baseline (0,2) 使整框落入界行格子不压版框线；
% 白字节点 x 不补偿、y +8.2pt（像素级墨迹实测 dx/dy≈0，L/R 留白约 20pt）。
\newcommand{\shuzanhead}[1]{%
  \begin{tikzpicture}[baseline={(0,2)}]
    \fill[sealbrown, rounded corners=6pt]
      (-76pt,-15pt) rectangle (76pt,15pt);
    \node[text=white, inner sep=0pt, anchor=center] at (0pt,8.2pt)
      {\shuzanfont\fontsize{23}{23}\selectfont #1};
  \end{tikzpicture}%
}
% ★ 修正书口卷名位置：chapter 框起始位置改回 cls 默认值，避免全局影响。
%   只在每个卷首页通过 \juan 命令里的空命令微调该页位置（安全，不会跨页污染）。
\ExplSyntaxOn
\AtBeginDocument{
  \dim_set:Nn \l__zhvt_shukou_chapter_title_start_dim
    { \l__zhvt_shukou_title_height_dim + \l__zhvt_font_size_dim * 2 }
  \dim_set:Nn \l__zhvt_shukou_page_number_start_dim
    { \l__zhvt_shukou_chapter_title_start_dim + \l__zhvt_shukou_chapter_title_height_dim }
}
\ExplSyntaxOff

\begin{document}

% ★★ 版心（书口）排布覆写：把书名放到鱼尾之下（2026-10-04）
%
%   cls 默认版心自上而下：书名（title 静态框，在鱼尾之上）→ 鱼尾 → 篇名（chapter 动态框）→ 页码
%   本模板改为：          鱼尾（位置不动）→ 书名 → 篇名 → 页码分割线 → 页码
%
%   手法：把 cls 放在 title 静态框里的书名，改放进 chapter 动态框（鱼尾之下）；
%         title 静态框留空。书名取 \l__zhvt_title_str，即 \title{...} 的第一参数，
%         所以换书时只改 \title，不必在此重复写书名。
%
%   ⚠ 三条必须遵守的写法（2026-10-04 实测）：
%     ① 汉字必须用 \hbox{...} 包住。flowfram 会把框内容装进一个宽 14.4mm（= 一个 line_sep
%        = 书口栏宽）的纵向 minipage 并按该宽度折行，而整页靠 PDF /Rotate 90 显示、
%        框内文字本身不变向 —— 不包 \hbox 时「围炉夜话」4 字会被逐字折成 4 行，
%        在版面上就成了横排 4 字并横向溢出版框。
%     ② 竖排下**一个 \hbox 就是一整列**。书名与篇名要同处一列，必须放进同一个 \hbox；
%        分成两个 \hbox 时，后一个会被当成新的一列、甩到隔壁栏去
%        （实测「序」跑到书口栏外侧 14.4mm 的正文列上）。中间的 \quad 在框内是横向空白，
%        显示出来正好是书名与篇名之间的纵向间隔。
%     ③ ★ 内容盒要写成 **与框等宽的定宽盒** `\hbox to \l__zhvt_shukou_chapter_title_height_dim{…\hfil}`。
%        实测：chapter 框是把内容盒沿纵向**居中**摆放的，内容宽度一变，书名就跟着漂：
%        自然宽度的 \hbox（书名 38.4mm）落在 x=157.2mm，加篇名后（57.6mm）落在 x=152.4mm，
%        目录页（无篇名）与正文页的书名就差半格。做成等宽盒后，无论有没有篇名，
%        内容固定从框顶起排，实测两页均落在 x=142.8mm，不再漂。
\ExplSyntaxOn
% 把书口 chapter 框（书名+篇名）起始位置下移 1 字高，避免 4 字卷名顶到鱼尾
% ⚠ cls 的 \__zhvt_calculate_dimensions: 在 \AtBeginDocument 里执行，会重算该 dim；
%   必须把覆写挂到 begin document 之后，否则会被二次覆盖（2026-10-05 实测）。
% (旧方案已移到导言区，直接修改 chapter 框起始位置)
% 卷标题命令 #1=书口篇名(鱼尾) #2=正文卷名（含卷序，卷名内空格用 ~）
%
% ★★ 目录页两项修订（2026-10-05，愚大人指出）★★
%   ① 页码起点统一对齐：
%      原状 —— 目录项写成 \par\hspace{4em}卷名\hspace{4em}页码，页码起点 = 卷名实宽 + 4em。
%      卷名实宽随字数变：「卷一~梁惠王上」6 字 +~、「卷七~离娄上」5 字 +~，
%      差 1 个字位 → 卷七~卷十 4 条的页码整体比其余 10 条高 27.21pt（一个字），看着上下不齐。
%      改法 —— 卷名补全角空格（宽 = 1 字 = 27.215pt）补足到 **7 字符位**再写进目录，
%      14 条页码起点遂全部落在同一条竖线上（起点对齐；末字随页码字数自然参差）。
%      ⚠ 只对目录里的副本 \l_tmpc_tl 补，正文卷名仍用 \l_tmpa_tl 原样 —— 正文像素不受影响。
%   ② 卷名整体上移半字位：前导 \hspace 由 4em 改 3.5em。
%      实测原状目录块在版框内偏下（天头留白 139.6pt / 地脚留白 110.7pt），
%      上移 0.5em = 13.6pt 后上下基本均分，视觉居中。
\NewDocumentCommand { \juan } { mm } {
	\clearpage
	\tl_set:Nn \l__zhvt_chapter_title_tl { #1 }
	\tl_set:Nn \l_tmpa_tl { #2 }
	\tl_set_eq:NN \l_tmpc_tl \l_tmpa_tl
	\int_step_inline:nn { 7 - \str_count:n { #2 } }
	  { \tl_put_right:Nn \l_tmpc_tl { 　 } }
	\tl_set:Nx \l_tmpb_tl { \zhnum { page } }
	\tl_put_left:Nn \l_tmpb_tl { \hspace { 4em } }
	\tl_put_left:Nn \l_tmpb_tl { \l_tmpc_tl }
	\tl_put_left:Nn \l_tmpb_tl { \par \hspace { 3.5em } }
	\addtocontents { toc } { \l_tmpb_tl }
	\skip_horizontal:n { \l__zhvt_font_size_dim * 4 }
	\l_tmpa_tl
	\par
	\thispagestyle { empty }
}
\cs_set_protected:Npn \__zhvt_set_shukou_content: {
	% ★ 鱼尾上方留白（title 静态框）嵌入闲叟钤印：
	%   整页 /Rotate 90，故此处把图用 \rotatebox{90} 转正为竖章；
	%   尺寸卡在版心条幅内（书内约 7mm 宽 × 33.8mm 高，留余量不压鱼尾）。
	\setstaticcontents * { title } {
		\vbox_to_ht:nn { \l__zhvt_baseline_skip_dim } {
			\vskip 14.2pt
			\hbox_to_wd:nn { \l__zhvt_shukou_title_height_dim } {
				\hfil
				\rotatebox [ origin = c ] { 90 } { \includegraphics [ width = 7mm , height = 33.8mm ] { yshs.png } }
				\hfil
			}
		}
	}
% ★ 版心（书口）统一改法（2026-11-16 愚大人定）：18pt＋框内居中＋书口文字同版框色
%   ① 原 27.2pt 时 7 字长篇名（如「廉颇蔺相如列传」）末字溢出 chapter 框、压住页码；
%      改 \fontsize{18pt}{18pt}，前后 \hfil 令字串在框内居中（不随篇名字数漂移）。
%   ② \vbox{\vskip 8.2pt} 定量下移，抵消缩字号后栏宽方向 8.2pt 上偏（像素实测 y 中心 543.3 = 原版）。
%      ⚠ \raisebox / \vspace* 在 flowfram 框内无效；换字号须同步换补偿：16pt→10.1、19pt→7.4。
%   ③ chapter 与 page~number 均置 \color{sealbrown}（HTML 874434），与版框／界行同色；正文仍黑。
%   ⚠ 页码字位、钤印、鱼尾、界行、正文一律未动；五部合订本同此口径。
	\setdynamiccontents * { chapter }
	{ \color{sealbrown} \vbox { \vskip 8.2pt \hbox to \l__zhvt_shukou_chapter_title_height_dim
		{ \hfil \fontsize{18pt}{18pt}\selectfont \l__zhvt_title_str \quad \l__zhvt_chapter_title_tl \hfil } } }
	\setdynamiccontents * { page~number } { \color{sealbrown} \quad \zhnum { page } }
}
\__zhvt_set_shukou_content:
\ExplSyntaxOff

\maketitle

% 单篇，无多卷目录，不渲染目录页

\juan{司马相如列传}{卷一百十七~司马相如列传第五十七}
\hspace*{2em}司马相如者。蜀郡成都人也。字长卿。少时好读书。学击剑。故其亲名之曰犬子。相如既学。慕蔺相如之为人。更名相如。以赀为郎。事孝景帝。为武骑常侍。非其好也。会景帝不好辞赋。是时梁孝王来朝。从游说之士齐人邹阳。淮阴枚乘。吴庄忌夫子之徒。相如见而说之。因病免。客游梁。梁孝王令与诸生同舍。相如得与诸生游士居数岁。乃著子虚之赋。\\
\hspace*{2em}会梁孝王卒。相如归。而家贫。无以自业。素与临邛令王吉相善。吉曰。﹁长卿久宦游不遂。而来过我。﹂于是相如往。舍都亭。临邛令缪为恭敬。日往朝相如。相如初尚见之。后称病。使从者谢吉。吉愈益谨肃。临邛中多富人。而卓王孙家僮八百人。程郑亦数百人。二人乃相谓曰。﹁令有贵客。为具召之。﹂并召令。令既至。卓氏客以百数。至日中。谒司马长卿。长卿谢病不能往。临邛令不敢尝食。自往迎相如。相如不得已。彊往。一坐尽倾。酒酣。临邛令前奏琴曰。﹁窃闻长卿好之。愿以自娱。﹂相如辞谢。为鼓一再行。是时卓王孙有女文君新寡。好音。故相如缪与令相重。而以琴心挑之。相如之临邛。从车骑。雍容间雅甚都。及饮卓氏。弄琴。文君窃从户窥之。心悦而好之。恐不得当也。既罢。相如乃使人重赐文君侍者通殷勤。文君夜亡奔相如。相如乃与驰归成都。家居徒四壁立。卓王孙大怒曰。﹁女至不材。我不忍杀。不分一钱也。﹂人或谓王孙。王孙终不听。文君久之不乐。曰。﹁长卿第俱如临邛。从昆弟假贷犹足为生。何至自苦如此。﹂相如与俱之临邛。尽卖其车骑。买一酒舍酤酒。而令文君当炉。相如身自著犊鼻裈。与保庸杂作。涤器于市中。卓王孙闻而耻之。为杜门不出。昆弟诸公更谓王孙曰。﹁有一男两女。所不足者非财也。今文君已失身于司马长卿。长卿故倦游。虽贫。其人材足依也。且又令客。独柰何相辱如此。﹂卓王孙不得已。分予文君僮百人。钱百万。及其嫁时衣被财物。文君乃与相如归成都。买田宅。为富人。\\
\hspace*{2em}居久之。蜀人杨得意为狗监。侍上。上读子虚赋而善之。曰。﹁朕独不得与此人同时哉。﹂得意曰。﹁臣邑人司马相如自言为此赋。﹂上惊。乃召问相如。相如曰。﹁有是。然此乃诸侯之事。未足观也。请为天子游猎赋。赋成奏之。﹂上许。令尚书给笔札。相如以﹁子虚﹂。虚言也。为楚称。﹁乌有先生﹂者。乌有此事也。为齐难。﹁无是公﹂者。无是人也。明天子之义。故空藉此三人为辞。以推天子诸侯之苑囿。其卒章归之于节俭。因以风谏。奏之天子。天子大说。其辞曰。\\
\hspace*{2em}楚使子虚使于齐。王悉发车骑。与使者出畋。畋罢。子虚过奼乌有先生。亡是公在焉。坐定。乌有先生问曰。﹁今日畋乐乎。﹂子虚曰。﹁乐。﹂﹁获多乎。﹂曰。﹁少。﹂﹁然则何乐。﹂对曰。﹁仆乐齐王之欲夸仆以车骑之众。而仆对以云梦之事也。﹂曰。﹁可得闻乎。﹂\\
\hspace*{2em}子虚曰。可。王车驾千乘。选徒万骑。畋于海滨。列卒满泽。罘罔弥山。掩兔辚鹿。射麋脚麟。骛于盐浦。割鲜染轮。射中获多。矜而自功。顾谓仆曰。﹃楚亦有平原广泽游猎之地饶乐若此者乎。楚王之猎孰与寡人乎。﹄仆下车对曰。﹃臣。楚国之鄙人也。幸得宿卫十有余年。时从出游。游于后园。览于有无。然犹未能遍睹也。又焉足以言其外泽者乎。﹄齐王曰。﹃虽然。略以子之所闻见而言之。﹄\\
\hspace*{2em}仆对曰。唯唯。臣闻楚有七泽。尝见其一。未睹其余也。臣之所见。盖特其小小耳者。名曰云梦。云梦者。方九百里。其中有山焉。其山则盘纡茀郁。隆崇嵂崒。岑崟参差。日月蔽亏。交错纠纷。上干青云。罢池陂陀。下属江河。其土则丹青赭垩。雌黄白坿。锡碧金银。众色炫耀。照烂龙鳞。其石则赤玉玫瑰。琳瑉琨吾。瑊玏玄厉。碝石碔玞。其东则有蕙圃。衡兰芷若。芎藭昌蒲。茳蓠麋芜。诸柘巴苴。其南则有平原广泽。登降陁靡。案衍坛曼。缘以大江。限以巫山。其高燥则生葴菥苞荔。薛莎青薠。其卑湿则生藏莨蒹葭。东蔷雕胡。莲藕觚卢。菴闾轩于。众物居之。不可胜图。其西则有涌泉清池。激水推移。外发芙蓉菱华。内隐钜石白沙。其中则有神龟蛟鼍。瑇瑁鳖鼋。其北则有阴林。其树楩柟豫章。桂椒木兰。檗离朱杨。樝梨梬栗。橘柚芬芳。其上则有鹓雏孔鸾。腾远射干。其下则有白虎玄豹。蟃蜒貙豻。兕象野犀。穷奇獌狿。\\
\hspace*{2em}于是乃使专诸之伦。手格此兽。楚王乃驾驯驳之驷。乘雕玉之舆。靡鱼须之桡旃。曳明月之珠旗。建干将之雄戟。左乌号之雕弓。右夏服之劲箭。阳子骖乘。纤阿为御。案节未舒。即陵狡兽。蹴蛩蛩。辚距虚。轶野马。轊陶駼。乘遗风。射游骐。倏眒倩浰。雷动猋至。星流霆击。弓不虚发。中必决眦。洞胸达腋。绝乎心系。获若雨兽。揜草蔽地。于是楚王乃弭节徘徊。翱翔容与。览乎阴林。观壮士之暴怒。与猛兽之恐惧。徼郄受诎。殚睹众物之变态。\\
\hspace*{2em}于是郑女曼姬。被阿锡。揄纻缟。櫜纤罗。垂雾縠。襞积褰绉。纡徐委曲。郁桡谿谷。衯衯裶裶。扬袘恤削。蜚纤垂髾。扶与猗靡。吸呷萃蔡。下摩兰蕙。上拂羽盖。错翡翠之威蕤。缪绕玉绥。缥乎忽忽。若神仙之仿佛。\\
\hspace*{2em}于是乃相与獠于蕙圃。媻珊勃窣上金隄。揜翡翠。射鵕璘。微矰出。纤缴施。弋白鹄。连驾鹅。双鸧下。玄鹤加。怠而后发。游于清池。浮文鹢。扬桂枻。张翠帷。建羽盖。罔玳瑁。钓紫贝。摐金鼓。吹鸣籁。榜人歌。声流喝。水蟲骇。波鸿沸。涌泉起。奔扬会。礧石相击。硠硠潏潏。若雷霆之声。闻乎数百里之外。\\
\hspace*{2em}﹁﹃将息獠者。击灵鼓。起烽燧。车案行。骑就队。纚乎淫淫。班乎裔裔。于是楚王乃登阳云之台。泊乎无为。澹乎自持。勺药之和具而后御之。不若大王终日驰骋而不下舆。脟割轮淬。自以为娱。臣窃观之。齐殆不如。﹄于是王默然无以应仆也。﹂\\
\hspace*{2em}乌有先生曰。﹁是何言之过也。足下不远千里。来况齐国。王悉发境内之士。而备车骑之众。以出田。乃欲戮力致获。以娱左右也。何名为夸哉。问楚地之有无者。原闻大国之风烈。先生之馀论也。今足下不称楚王之德厚。而盛推云梦以为高。奢言淫乐而显侈靡。窃为足下不取也。必若所言。固非楚国之美也。有而言之。是章君之恶。无而言之。是害足下之信。章君之恶而伤私义。二者无一可。而先生行之。必且轻于齐而累于楚矣。且齐东陼巨海。南有琅邪。观乎成山。射乎之罘。浮勃澥。游孟诸。邪与肃慎为邻。右以汤谷为界。秋田乎青丘。傍徨乎海外。吞若云梦者八九。其于胸中曾不蒂芥。若乃俶傥瑰伟。异方殊类。珍怪鸟兽。万端鳞萃。充仞其中者。不可胜记。禹不能名。契不能计。然在诸侯之位。不敢言游戏之乐。苑囿之大。先生又见客。是以王辞而不复。何为无用应哉。﹂\\
\hspace*{2em}无是公听然而笑曰。楚则失矣。齐亦未为得也。夫使诸侯纳贡者。非为财币。所以述职也。封疆画界者。非为守御。所以禁淫也。今齐列为东籓。而外私肃慎。捐国逾限。越海而田。其于义故未可也。且二君之论。不务明君臣之义而正诸侯之礼。徒事争游猎之乐。苑囿之大。欲以奢侈相胜。荒淫相越。此不可以扬名发誉。而適足以贬君自损也。且夫齐楚之事又焉足道邪。君未睹夫巨丽也。独不闻天子之上林乎。\\
\hspace*{2em}左苍梧。右西极。丹水更其南。紫渊径其北。终始灞浐。出入泾渭。酆镐潦潏。纡馀委蛇。经营乎其内。荡荡乎八川分流。相背而异态。东西南北。驰骛往来。出乎椒丘之阙。行乎洲淤之浦。经乎桂林之中。过乎泱漭之野。汩乎混流。顺阿而下。赴隘狭之口。触穹石。激堆埼。沸乎暴怒。汹涌澎湃。滭弗宓汩。逼侧泌瀄。横流逆折。转腾潎冽。滂濞沆溉。穹隆云桡。宛潬胶盭。逾波趋浥。涖涖下濑。批岩冲拥。奔扬滞沛。临坻注壑。瀺灂霣坠。沈沈隐隐。砰磅訇礚。潏潏淈淈。湁潗鼎沸。驰波跳沫。汩濦漂疾。悠远长怀。寂漻无声。肆乎永归。然后灏溔潢漾。安翔徐回。翯乎滈滈。东注太湖。衍溢陂池。于是乎鲛龙赤螭。䱭䲛渐离。鰅鳙鳍鮀。禺禺魼鳎。揵鳍掉尾。振鳞奋翼。潜处乎深岩。鱼鳖讙声。万物众伙。明月珠子。的皪江靡。蜀石黄碝。水玉磊砢。磷磷烂烂。采色澔汗。藂积乎其中。鸿鹔鹄鸨。鴐鹅属玉。交精旋目。烦鹜庸渠。箴疵䴔卢。群浮乎其上。泛淫泛滥。随风澹淡。与波摇荡。奄薄水渚。唼喋菁藻。咀嚼菱藕。\\
\hspace*{2em}于是乎崇山矗矗。巃嵷崔巍。深林巨木。崭岩㟥嵳。九嵕嶻嶭。南山峨峨。岩陁甗崎。摧崣崛崎。振溪通谷。蹇产沟渎。谽呀豁閕。阜陵别岛。崴磈葨廆。丘虚堀礨。隐辚郁垒。登降陁靡。陂池貏豸。沇溶淫鬻。散涣夷陆。亭皋千里。靡不被筑。揜以绿蕙。被以江蓠。糅以蘪芜。杂以留夷。布结缕。攒戾莎。揭车衡兰。槀本射干。茈姜蘘荷。葴持若荪。鲜支黄砾。蒋苎青薠。布濩闳泽。延曼太原。离靡广衍。应风披靡。吐芳扬烈。郁郁菲菲。众香发越。肸蚃布写。晻薆咇茀。于是乎周览泛观。缜纷轧芴。芒芒恍忽。视之无端。察之无涯。日出东沼。入乎西陂。其南则隆冬生长。涌水跃波。其兽则㺎旄貘犛。沈牛麈麋。赤首圜题。穷奇象犀。其北则盛夏含冻裂地。涉冰揭河。其兽则麒麟角端。騊駼橐驼。蛩蛩驒騱。駃騠驴骡。\\
\hspace*{2em}于是乎离宫别馆。弥山跨谷。高廊四注。重坐曲阁。华榱璧珰。辇道纚属。步櫩周流。长途中宿。夷嵕筑堂。累台增成。岩窔洞房。頫杳眇而无见。仰攀橑而扪天。奔星更于闺闼。宛虹扡于楯轩。青龙蚴蟉于东箱。象舆婉僤于西清。灵圄燕于闲馆。偓佺之伦。暴于南荣。醴泉涌于清室。通川过于中庭。盘石振崖。嵚岩倚倾。嵯峨磼礏。刻削峥嵘。玫瑰碧琳。珊瑚丛生。琘玉旁唐。玢豳文鳞。赤瑕驳荦。杂臿其间。晁采琬琰。和氏出焉。\\
\hspace*{2em}于是乎卢橘夏熟。黄甘橙楱。枇杷橪柿。亭奈厚朴。梬枣杨梅。樱桃蒲陶。隐夫薁棣。荅遝离支。罗乎后宫。列乎北园。貤丘陵。下平原。扬翠叶。扤紫茎。发红华。垂朱荣。煌煌扈扈。照曜钜野。沙棠栎槠。华枫枰栌。留落胥邪。仁频并闾。欃檀木兰。豫章女贞。长千仞。大连抱。夸条直畅。实叶葰楙。攒立丛倚。连卷欐佹。崔错癹骫。坑衡閜砢。垂条扶疏。落英幡纚。纷溶箾蔘。猗狔从风。藰莅卉歙。盖象金石之声。管籥之音。偨池茈虒。旋还乎后宫。杂袭累辑。被山缘谷。循阪下隰。视之无端。究之无穷。\\
\hspace*{2em}于是乎玄猨素雌。蜼玃飞鸓。蛭蜩蠼猱。獑胡豰蛫。栖息乎其间。长啸哀鸣。翩幡互经。夭蟜枝格。偃蹇杪颠。隃绝梁。腾殊榛。捷垂条。掉希间。牢落陆离。烂漫远迁。\\
\hspace*{2em}若此者数百千处。娱游往来。宫宿馆舍。庖厨不徙。后宫不移。百官备具。\\
\hspace*{2em}﹁于是乎背秋涉冬。天子校猎。乘镂象。六玉虬。拖蜺旌。靡云旗。前皮轩。后道游。孙叔奉辔。卫公参乘。扈从横行。出乎四校之中。鼓严簿。纵猎者。河江为阹。泰山为橹。车骑雷起。殷天动地。先后陆离。离散别追。淫淫裔裔。缘陵流泽。云布雨施。﹂\\
\hspace*{2em}生貔豹。搏豺狼。手熊罴。足壄羊。蒙鹖苏。绔白虎。被班文。跨壄马。凌三嵕之危。下碛历之坻。径峻赴险。越壑厉水。椎蜚廉。弄獬豸。格虾蛤。鋋猛氏。羂騕褭。射封豕。箭不苟害。解脰陷脑。弓不虚发。应声而倒。于是乘舆弭节徘徊。翱翔往来。睨部曲之进退。览将帅之变态。然后侵淫促节。儵夐远去。流离轻禽。蹴履狡兽。轊白鹿。捷狡兔。轶赤电。遗光耀。追怪物。出宇宙。弯蕃弱。满白羽。射游枭。栎蜚遽。择肉而后发。先中而命处。弦矢分。艺殪仆。\\
\hspace*{2em}然后扬节而上浮。陵惊风。历骇梠。乘虚无。与神俱。辚玄鹤。乱昆鸡。遒孔鸾。促鵕璘。拂鹥鸟。捎凤皇。捷鸳雏。掩焦明。\\
\hspace*{2em}道尽途殚。回车而还。消遥乎襄羊。降集乎北纮。率乎直指。晻乎反乡。蹷石阙。历封峦。过鳷鹊。望露寒。下棠梨。息宜春。西驰宣曲。濯鹢牛首。登龙台。掩细柳。观士大夫之勤略。均猎者之所得获。徒车之所轥轹。步骑之所蹂若。人臣之所蹈籍。与其穷极倦谻。惊惮詟伏。不被创刃而死者。他他籍籍。填坑满谷。掩平弥泽。\\
\hspace*{2em}于是乎游戏懈怠。置酒乎颢天之台。张乐乎轇輵之宇。撞千石之钟。立万石之虡。建翠华之旗。树灵鼍之鼓。奏陶唐氏之舞。听葛天氏之歌。千人唱。万人和。山陵为之震动。川谷为之荡波。巴渝宋蔡。淮南干遮。文成颠歌。族居递奏。金鼓迭起。铿鎗闛鞈。洞心骇耳。荆吴郑卫之声。韶濩武象之乐。阴淫案衍之音。鄢郢缤纷。激楚结风。俳优侏儒。狄鞮之倡。所以娱耳目乐心意者。丽靡烂漫于前。靡曼美色于后。\\
\hspace*{2em}若夫青琴。宓妃之徒。绝殊离俗。妖冶娴都。靓妆刻饰。便嬛绰约。柔桡嫚嫚。妩媚纤弱。曳独茧之褕绁。眇阎易以恤削。便姗嫳屑。与俗殊服。芬芳沤郁。酷烈淑郁。皓齿粲烂。宜笑的皪。长眉连娟。微睇绵藐。色授魂与。心愉于侧。\\
\hspace*{2em}于是酒中乐酣。天子芒然而思。似若有亡。曰。嗟乎。此大奢侈。朕以览听馀闲。无事弃日。顺天道以杀伐。时休息于此。恐后叶靡丽。遂往而不返。非所以为继嗣创业垂统也。于是乎乃解酒罢猎。而命有司曰。﹃地可垦辟。悉为农郊。以赡萌隶。隤墙填堑。使山泽之人得至焉。实陂池而勿禁。虚宫馆而勿仞。发仓廪以救贫穷。补不足。恤鳏寡。存孤独。出德号。省刑罚。改制度。易服色。革正朔。与天下为更始。﹄\\
\hspace*{2em}于是历吉日以斋戒。袭朝服。乘法驾。建华旗。鸣玉鸾。游于六艺之囿。驰骛乎仁义之涂。览观春秋之林。射狸首。兼驺虞。弋玄鹤。舞干戚。载云䍐。揜群雅。悲伐檀。乐乐胥。修容乎礼园。翱翔乎书圃。述易道。放怪兽。登明堂。坐清庙。次群臣。奏得失。四海之内。靡不受获。于斯之时。天下大说。乡风而听。随流而化。卉然兴道而迁义。刑错而不用。德隆于三王。而功羡于五帝。若此故猎。乃可喜也。\\
\hspace*{2em}﹁若夫终日驰骋。劳神苦形。罢车马之用。抏士卒之精。费府库之财。而无德厚之恩。务在独乐。不顾众庶。亡国家之政。贪雉兔之获。则仁者不繇也。从此观之。齐楚之事。岂不哀哉。地方不过千里。而囿居九百。是草木不得垦辟。而人无所食也。夫以诸侯之细。而乐万乘之侈。仆恐百姓被其尤也。﹂\\
\hspace*{2em}于是二子愀然改容。超若自失。逡巡避席。曰。﹁鄙人固陋。不知忌讳。乃今见教。谨受命矣。﹂\\
\hspace*{2em}赋奏。天子以为郎。无是公言天子上林广大。山谷水泉万物。乃子虚言楚云梦所有甚众。侈靡过其实。且非义理所尚。故删取其要。归正道而论之。\\
\hspace*{2em}相如为郎数岁。会唐蒙使略通夜郎西僰中。发巴蜀吏卒千人。郡又多为发转漕万馀人。用兴法诛其渠帅。巴蜀民大惊恐。上闻之。乃使相如责唐蒙。因喻告巴蜀民以非上意。檄曰。\\
\hspace*{2em}告巴蜀太守。蛮夷自擅不讨之日久矣。时侵犯边境。劳士大夫。陛下即位。存抚天下。辑安中国。然后兴师出兵。北征匈奴。单于怖骇。交臂受事。诎膝请和。康居西域。重译请朝。稽首来享。移师东指。闽越相诛。右吊番禺。太子入朝。南夷之君。西僰之长。常效贡职。不敢怠堕。延颈举踵。喁喁然皆争归义。欲为臣妾。道里辽远。山川阻深。不能自致。夫不顺者已诛。而为善者未赏。故遣中郎将往宾之。发巴蜀士民各五百人。以奉币帛。卫使者不然。靡有兵革之事。战斗之患。今闻其乃发军兴制。惊惧子弟。忧患长老。郡又擅为转粟运输。皆非陛下之意也。当行者或亡逃自贼杀。亦非人臣之节也。\\
\hspace*{2em}夫边郡之士。闻烽举燧燔。皆摄弓而驰。荷兵而走。流汗相属。唯恐居后。触白刃。冒流矢。义不反顾。计不旋踵。人怀怒心。如报私雠。彼岂乐死恶生。非编列之民。而与巴蜀异主哉。计深虑远。急国家之难。而乐尽人臣之道也。故有剖符之封。析珪而爵。位为通侯。居列东第。终则遗显号于后世。传土地于子孙。行事甚忠敬。居位甚安佚。名声施于无穷。功烈著而不灭。是以贤人君子。肝脑涂中原。膏液润野草而不辞也。今奉币役至南夷。即自贼杀。或亡逃抵诛。身死无名。谥为至愚。耻及父母。为天下笑。人之度量相越。岂不远哉。然此非独行者之罪也。父兄之教不先。子弟之率不谨也。寡廉鲜耻。而俗不长厚也。其被刑戮。不亦宜乎。\\
\hspace*{2em}陛下患使者有司之若彼。悼不肖愚民之如此。故遣信使晓喻百姓以发卒之事。因数之以不忠死亡之罪。让三老孝弟以不教诲之过。方今田时。重烦百姓。已亲见近县。恐远所谿谷山泽之民不遍闻。檄到。亟下县道。使咸知陛下之意。唯毋忽也。\\
\hspace*{2em}相如还报。唐蒙已略通夜郎。因通西南夷道。发巴。蜀。广汉卒。作者数万人。治道二岁。道不成。士卒多物故。费以巨万计。蜀民及汉用事者多言其不便。是时邛筰之君长闻南夷与汉通。得赏赐多。多欲原为内臣妾。请吏。比南夷。天子问相如。相如曰。﹁邛。筰。厓。駹者近蜀。道亦易通。秦时尝通为郡县。至汉兴而罢。今诚复通。为置郡县。愈于南夷。﹂天子以为然。乃拜相如为中郎将。建节往使。副使王然于。壶充国。吕越人驰四乘之传。因巴蜀吏币物以赂西夷。至蜀。蜀太守以下郊迎。县令负弩矢先驱。蜀人以为宠。于是卓王孙。临邛诸公皆因门下献牛酒以交驩。卓王孙喟然而叹。自以得使女尚司马长卿晚。而厚分与其女财。与男等同。司马长卿便略定西夷。邛。筰。厓。駹。斯榆之君皆请为内臣。除边关。关益斥。西至沬。若水。南至牂柯为徼。通零关道。桥孙水以通邛都。还报天子。天子大说。\\
\hspace*{2em}相如使时。蜀长老多言通西南夷不为用。唯大臣亦以为然。相如欲谏。业已建之。不敢。乃著书。籍以蜀父老为辞。而己诘难之。以风天子。且因宣其使指。令百姓知天子之意。其辞曰。\\
\hspace*{2em}汉兴七十有八载。德茂存乎六世。威武纷纭。湛恩汪濊。群生澍濡。洋溢乎方外。于是乃命使西征。随流而攘。风之所被。罔不披靡。因朝厓从駹。定筰存邛。略斯榆。举苞满。结轶还辕。东乡将报。至于蜀都。\\
\hspace*{2em}耆老大夫荐绅先生之徒二十有七人。俨然造焉。辞毕。因进曰。﹁盖闻天子之于夷狄也。其义羁縻勿绝而已。今罢三郡之士。通夜郎之涂。三年于兹。而功不竟。士卒劳倦。万民不赡。今又接以西夷。百姓力屈。恐不能卒业。此亦使者之累也。窃为左右患之。且夫邛。筰。西僰之与中国并也。历年兹多。不可记已。仁者不以德来。彊者不以力并。意者其殆不可乎。今割齐民以附夷狄。弊所恃以事无用。鄙人固陋。不识所谓。﹂\\
\hspace*{2em}使者曰。乌谓此邪。必若所云。则是蜀不变服而巴不化俗也。余尚恶闻若说。然斯事体大。固非观者之所觏也。余之行急。其详不可得闻已。请为大夫粗陈其略。\\
\hspace*{2em}盖世必有非常之人。然后有非常之事。有非常之事。然后有非常之功。非常者。固常之所异也。故曰非常之原。黎民惧焉。及臻厥成。天下晏如也。\\
\hspace*{2em}昔者鸿水浡出。氾滥衍溢。民人登降移徙。陭麕而不安。夏后氏戚之。乃堙鸿水。决江疏河。漉沈赡菑。东归之于海。而天下永宁。当斯之勤。岂唯民哉。心烦于虑而身亲其劳。躬胝无胈。肤不生毛。故休烈显乎无穷。声称浃乎于兹。\\
\hspace*{2em}且夫贤君之践位也。岂特委琐握麀。拘文牵俗。循诵习传。当世取说云尔哉。必将崇论闳议。创业垂统。为万世规。故驰骛乎兼容并包。而勤思乎参天贰地。且诗不云乎。﹃普天之下。莫非王土。率土之滨。莫非王臣。﹄是以六合之内。八方之外。浸浔衍溢。怀生之物有不浸润于泽者。贤君耻之。今封疆之内。冠带之伦。咸获嘉祉。靡有阙遗矣。而夷狄殊俗之国。辽绝异党之地。舟舆不通。人迹罕至。政教未加。流风犹微。内之则犯义侵礼于边境。外之则邪行横作。放弑其上。君臣易位。尊卑失序。父兄不辜。幼孤为奴。系累号泣。内乡而怨。曰﹃盖闻中国有至仁焉。德洋而恩普。物靡不得其所。今独曷为遗己﹄。举踵思慕。若枯旱之望雨。盭夫为之垂涕。况乎上圣。又恶能已。故北出师以讨彊胡。南驰使以诮劲越。四面风德。二方之君鳞集仰流。原得受号者以亿计。故乃关沬。若。徼牂柯。镂零山。梁孙原。创道德之涂。垂仁义之统。将博恩广施。远抚长驾。使疏逖不闭。阻深闇昧得耀乎光明。以偃甲兵于此。而息诛伐于彼。遐迩一体。中外提福。不亦康乎。夫拯民于沈溺。奉至尊之休德。反衰世之陵迟。继周氏之绝业。斯乃天子之急务也。百姓虽劳。又恶可以已哉。\\
\hspace*{2em}﹁且夫王事固未有不始于忧勤。而终于佚乐者也。然则受命之符。合在于此矣。方将增泰山之封。加梁父之事。鸣和鸾。扬乐颂。上咸五。下登三。观者未睹指。听者未闻音。犹鹪明已翔乎寥廓。而罗者犹视乎薮泽。悲夫。﹂\\
\hspace*{2em}于是诸大夫芒然丧其所怀来而失厥所以进。喟然并称曰。﹁允哉汉德。此鄙人之所原闻也。百姓虽怠。请以身先之。﹂敞罔靡徙。因迁延而辞避。\\
\hspace*{2em}其后人有上书言相如使时受金。失官。居岁馀。复召为郎。\\
\hspace*{2em}相如口吃而善著书。常有消渴疾。与卓氏婚。饶于财。其进仕宦。未尝肯与公卿国家之事。称病间居。不慕官爵。常从上至长杨猎。是时天子方好自击熊彘。驰逐野兽。相如上疏谏之。其辞曰。臣闻物有同类而殊能者。故力称乌获。捷言庆忌。勇期贲。育。臣之愚。窃以为人诚有之。兽亦宜然。今陛下好陵阻险。射猛兽。卒然遇轶材之兽。骇不存之地。犯属车之清尘。舆不及还辕。人不暇施巧。虽有乌获。逢蒙之伎。力不得用。枯木朽株尽为害矣。是胡越起于毂下。而羌夷接轸也。岂不殆哉。虽万全无患。然本非天子之所宜近也。\\
\hspace*{2em}且夫清道而后行。中路而后驰。犹时有衔橛之变。而况涉乎蓬蒿。驰乎丘坟。前有利兽之乐而内无存变之意。其为祸也不亦难矣。夫轻万乘之重不以为安。而乐出于万有一危之涂以为娱。臣窃为陛下不取也。\\
\hspace*{2em}盖明者远见于未萌而智者避危于无形。祸固多藏于隐微而发于人之所忽者也。故鄙谚曰﹁家累千金。坐不垂堂﹂。此言虽小。可以喻大。臣原陛下之留意幸察。\\
\hspace*{2em}上善之。还过宜春宫。相如奏赋以哀二世行失也。其辞曰。\\
\hspace*{2em}登陂阤之长阪兮。坌入曾宫之嵯峨。临曲江之隑州兮。望南山之参差。岩岩深山之谾谾兮。通谷魌兮谽谺。汩淢噏习以永逝兮。注平皋之广衍。观众树之蓊薆兮。览之榛榛。东驰土山兮。北揭石濑。弥节容与兮。历吊二世。持身不谨兮。亡国失埶。信谗不寤兮。宗庙灭绝。呜呼哀哉。操行之不得兮。坟墓芜秽而不脩兮。魂无归而不食。夐邈绝而不齐兮。弥久远而愈鬐。精罔阆而飞扬兮。拾九天而永逝。呜呼哀哉。\\
\hspace*{2em}相如拜为孝文园令。天子既美子虚之事。相如见上好仙道。因曰。﹁上林之事未足美也。尚有靡者。臣尝为大人赋。未就。请具而奏之。﹂相如以为列仙之传居山泽间。形容甚癯。此非帝王之仙意也。乃遂就大人赋。其辞曰。\\
\hspace*{2em}世有大人兮。在于中州。宅弥万里兮。曾不足以少留。悲世俗之迫隘兮。朅轻举而远游。垂绛幡之素蜺兮。载云气而上浮。建格泽之长竿兮。总光耀之采旄。垂旬始以为幓兮。抴彗星而为髾。掉指桥以偃蹇兮。又旖旎以招摇。揽欃枪以为旌兮。靡屈虹而为绸。红杳渺以眩湣兮。猋风涌而云浮。驾应龙象舆之蠖略逶丽兮。骖赤螭青虬之鞮蟉蜿蜒。低卬夭蟜据以骄骜兮。诎折隆穷蠼以连卷沛艾赳螑仡以佁儗兮。放散畔岸骧以孱颜。跮踱輵辖容以委丽兮。绸缪偃蹇怵鞨以梁倚。纠蓼叫奡蹋以艐路兮。蔑蒙踊跃腾而狂趡。莅飒卉翕熛至电过兮。焕然雾除。霍然云消。\\
\hspace*{2em}邪绝少阳而登太阴兮。与真人乎相求。互折窈窕以右转兮。横厉飞泉以正东。悉徵灵圉而选之兮。部乘众神于瑶光。使五帝先导兮。反太一而从陵阳。左玄冥而右含雷兮。前陆离而后潏湟。厮征伯侨而役羡门兮。属岐伯使尚方。祝融惊而跸御兮。清雰气而后行。屯余车其万乘兮。綷云盖而树华旗。使句芒其将行兮。吾欲往乎南嬉。\\
\hspace*{2em}历唐尧于崇山兮。过虞舜于九疑。纷湛湛其差错兮。杂遝胶葛以方驰。骚扰冲苁其相纷挐兮。滂濞泱轧洒以林离。钻罗列聚丛以茏茸兮。衍曼流烂坛以陆离。径入雷室之砰磷郁律兮。洞出鬼谷之嚬礨嵬靺。遍览八纮而观四荒兮。朅渡九江而越五河。经营炎火而浮弱水兮。杭绝浮渚而涉流沙。奄息总极氾滥水嬉兮。使灵娲鼓瑟而舞冯夷。时若薆薆将混浊兮。召屏翳诛风伯而刑雨师。西望昆仑之轧沕洸忽兮。直径驰乎三危。排阊阖而入帝宫兮。载玉女而与之归。舒阆风而摇集兮。亢乌腾而一止。低回阴山翔以纡曲兮。吾乃今目睹西王母鱇然白首。载胜而穴处兮。亦幸有三足乌为之使。必长生若此而不死兮。虽济万世不足以喜。\\
\hspace*{2em}回车朅来兮。绝道不周。会食幽都。呼吸沆瀣餐朝霞。噍咀芝英兮叽琼华。僸侵浔而高纵兮。纷鸿涌而上厉。贯列缺之倒景兮。涉丰隆之滂沛。驰游道而脩降兮。骛遗雾而远逝。迫区中之隘陕兮。舒节出乎北垠。遗屯骑于玄阙兮。轶先驱于寒门。下峥嵘而无地兮。上寥廓而无天。视眩眠而无见兮。听惝恍而无闻。乘虚无而上假兮。超无友而独存。\\
\hspace*{2em}相如既奏大人之颂。天子大说。飘飘有凌云之气。似游天地之间意。\\
\hspace*{2em}相如既病免。家居茂陵。天子曰。﹁司马相如病甚。可往从悉取其书。若不然。后失之矣。﹂使所忠往。而相如已死。家无书。问其妻。对曰。﹁长卿固未尝有书也。时时著书。人又取去。即空居。长卿未死时。为一卷书。曰有使者来求书。奏之。无他书。﹂其遗札书言封禅事。奏所忠。忠奏其书。天子异之。其书曰。伊上古之初肇。自昊穹兮生民。历撰列辟。以迄于秦。率迩者踵武。逖听者风声。纷纶葳蕤。堙灭而不称者。不可胜数也。续昭夏。崇号谥。略可道者七十有二君。罔若淑而不昌。畴逆失而能存。\\
\hspace*{2em}轩辕之前。遐哉邈乎。其详不可得闻也。五三六经载籍之传。维见可观也。书曰﹁元首明哉。股肱良哉﹂。因斯以谈。君莫盛于唐尧。臣莫贤于后稷。后稷创业于唐。公刘发迹于西戎。文王改制。爰周郅隆。大行越成。而后陵夷衰微。千载无声。岂不善始善终哉。然无异端。慎所由于前。谨遗教于后耳。故轨迹夷易。易遵也。湛恩濛涌。易丰也。宪度著明。易则也。垂统理顺。易继也。是以业隆于繦褓而崇冠于二后。揆厥所元。终都攸卒。未有殊尤绝迹可考于今者也。然犹蹑梁父。登泰山。建显号。施尊名。大汉之德。逢涌原泉。沕潏漫衍。旁魄四塞。云尃雾散。上暢九垓。下溯八埏。怀生之类霑濡浸润。协气横流。武节飘逝。迩陕游原。迥阔泳沫。首恶湮没。闇昧昭晢。昆蟲凯泽。回首面内。然后囿驺虞之珍群。徼麋鹿之怪兽。鳒一茎六穗于庖。牺双觡共抵之兽。获周馀珍收龟于岐。招翠黄乘龙于沼。鬼神接灵圉。宾于间馆。奇物谲诡。俶傥穷变。钦哉。符瑞臻兹。犹以为薄。不敢道封禅。盖周跃鱼陨杭。休之以燎。微夫斯之为符也。以登介丘。不亦恧乎。进让之道。其何爽与。\\
\hspace*{2em}于是大司马进曰。﹁陛下仁育群生。义征不憓。诸夏乐贡。百蛮执贽。德侔往初。功无与二。休烈浃洽。符瑞众变。期应绍至。不特创见。意者泰山。梁父设坛场望幸。盖号以况荣。上帝垂恩储祉。将以荐成。陛下谦让而弗发也。挈三神之驩。缺王道之仪。群臣恧焉。或谓且天为质闇。珍符固不可辞。若然辞之。是泰山靡记而梁父靡几也。亦各并时而荣。咸济世而屈。说者尚何称于后。而云七十二君乎。夫修德以锡符。奉符以行事。不为进越。故圣王弗替。而修礼地祗。谒款天神。勒功中岳。以彰至尊。舒盛德。发号荣。受厚福。以浸黎民也。皇皇哉斯事。天下之壮观。王者之丕业。不可贬也。原陛下全之。而后因杂荐绅先生之略术。使获燿日月之末光绝炎。以展采错事。犹兼正列其义。校饬厥文。作春秋一艺。将袭旧六为七。摅之无穷。俾万世得激清流。扬微波。蜚英声。腾茂实。前圣之所以永保鸿名而常为称首者用此。宜命掌故悉奏其义而览焉。﹂\\
\hspace*{2em}于是天子沛然改容。曰。﹁愉乎。朕其试哉。﹂乃迁思回虑。总公卿之议。询封禅之事。诗大泽之博。广符瑞之富。乃作颂曰。\\
\hspace*{2em}自我天覆。云之油油。甘露时雨。厥壤可游。滋液渗漉。何生不育。嘉穀六穗。我穑曷蓄。\\
\hspace*{2em}非唯雨之。又润泽之。非唯濡之。氾尃濩之。万物熙熙。怀而慕思。名山显位。望君之来。君乎君乎。侯不迈哉。\\
\hspace*{2em}般般之兽。乐我君囿。白质黑章。其仪可。旼旼睦睦。君子之能。盖闻其声。今观其来。厥涂靡踪。天瑞之徵。兹亦于舜。虞氏以兴。\\
\hspace*{2em}濯濯之麟。游彼灵畤。孟冬十月。君俎郊祀。驰我君舆。帝以享祉。三代之前。盖未尝有。\\
\hspace*{2em}宛宛黄龙。兴德而升。采色炫燿。熿炳煇煌。正阳显见。觉寤黎烝。于传载之。云受命所乘。\\
\hspace*{2em}厥之有章。不必谆谆。依类讬寓。谕以封峦。\\
\hspace*{2em}披艺观之。天人之际已交。上下相发允答。圣王之德。兢兢翼翼也。故曰﹁兴必虑衰。安必思危﹂。是以汤武至尊严。不失肃祗。舜在假典。顾省厥遗。此之谓也。\\
\hspace*{2em}司马相如既卒五岁。天子始祭后土。八年而遂先礼中岳。封于太山。至梁父禅肃然。\\
\hspace*{2em}相如他所著。若遗平陵侯书。与五公子相难。草木书篇不采。采其尤著公卿者云。\\
\hspace*{2em}太史公曰。春秋推见至隐。易本隐之以显。大雅言王公大人而德逮黎庶。小雅讥小己之得失。其流及上。所以言虽外殊。其合德一也。相如虽多虚辞滥说。然其要归引之节俭。此与诗之风谏何异。杨雄以为靡丽之赋。劝百风一。犹驰骋郑卫之声。曲终而奏雅。不已亏乎。余采其语可论者著于篇。\\

% ===== 索隐述赞（司马贞《史记索隐》）· 统一样式 =====
% 题头「索隐述赞曰」五字并入反白框；正文每四句一列，定宽盒 \hfil 居列中；
% \fontsize{24}{41} 列距对齐全书界行（121+41k）；\clearpage 使述赞独占一页。
\clearpage
\shuzanhead{索隐述赞曰}\linebreak
{\shuzanbody\fontsize{24}{41}\selectfont
\hbox to \hsize{\hfil 相如纵诞。窃赀卓氏。其学无方。其才足倚。\hfil}\linebreak
\hbox to \hsize{\hfil 子虚过吒。上林非侈。四马还邛。百金献伎。\hfil}\linebreak
\hbox to \hsize{\hfil 惜哉封禅。遗文卓尔。\hfil}%
\par}
\end{document}
